\documentclass[a4paper]{article}
% Español (México) con XeLaTeX: fontspec para fuentes Unicode, polyglossia
% para la división silábica y los nombres del documento en la variante mexicana.
\usepackage{fontspec}
\usepackage{polyglossia}
\setdefaultlanguage[variant=mexican]{spanish}
% Los nombres chinos van como «pinyin (original)»: xeCJK da a los caracteres
% Han su propia fuente; sin ella XeLaTeX los omite con un aviso.
% FandolSong no cubre algunos caracteres de nombres (U+73FA); WenQuanYi Zen Hei sí.
\usepackage[AutoFallBack=true]{xeCJK}
\setCJKmainfont{FandolSong-Regular.otf}
\setCJKfallbackfamilyfont{\CJKrmdefault}{WenQuanYi Zen Hei}
% polyglossia no traduce los nombres de listings.
\AtBeginDocument{\providecommand{\lstlistingname}{}\renewcommand{\lstlistingname}{Listado}%
  \providecommand{\lstlistlistingname}{}\renewcommand{\lstlistlistingname}{Índice de listados}}
\usepackage{amsmath, amssymb}
\usepackage{graphicx}
\usepackage[margin=2.35cm]{geometry}
\usepackage[most]{tcolorbox}
\usepackage{booktabs}
\usepackage{tabularx}
\usepackage{longtable}
\usepackage{array}
\usepackage{float}
\usepackage{xcolor}
\usepackage{hyperref}
\usepackage{fancyhdr}
\setCJKmainfont[AutoFakeBold=2.4]{AR PL UMing CN}
\setCJKsansfont[AutoFakeBold=2.4]{Droid Sans Fallback}
\setCJKmonofont[AutoFakeBold=2.4]{Droid Sans Fallback}
\hypersetup{colorlinks=true, linkcolor=blue!55!black, urlcolor=blue!60!black}
\graphicspath{{./}{figures/}}
\setlength{\parskip}{0.55em}
\setlength{\parindent}{2em}
\setlength{\emergencystretch}{3em}
\renewcommand{\arraystretch}{1.18}
\newtcolorbox{knowledgebox}[1]{enhanced, colback=blue!5!white, colframe=blue!70!black, colbacktitle=blue!70!black, coltitle=white, fonttitle=\bfseries, title={#1}, sharp corners, breakable}
\newtcolorbox{importantbox}[1]{enhanced, colback=yellow!10!white, colframe=yellow!75!black, colbacktitle=yellow!75!black, coltitle=black, fonttitle=\bfseries, title={#1}, sharp corners, breakable}
\newtcolorbox{warningbox}[1]{enhanced, colback=red!5!white, colframe=red!70!black, colbacktitle=red!70!black, coltitle=white, fonttitle=\bfseries, title={#1}, sharp corners, breakable}
\pagestyle{fancy}
\fancyhf{}
\fancyfoot[C]{\thepage}
\fancyfoot[R]{\footnotesize\color{gray}\href{https://github.com/hqhq1025/ai-course-notes}{AI Course Notes}}
\renewcommand{\headrulewidth}{0pt}
\renewcommand{\footrulewidth}{0pt}
\newcommand{\videourl}{https://www.youtube.com/watch?v=qZbzFZ2R_Nw}
\newcommand{\repourl}{https://github.com/hqhq1025/ai-course-notes}

\begin{document}
\begin{titlepage}
\centering
{\Large Notas didácticas de una entrevista a fondo\par}
\vspace{1.0cm}
{\huge\bfseries Sistemas generativos, plataformas de contenido y «sin intermediarios que se queden con el margen»\par}
\vspace{0.5cm}
{\Large Zhang Xiaojun conversa con Wang Guan (王冠), fundador de ONE2X\par}
\vspace{0.35cm}
{\large Elaborado a partir del video público de Zhang Xiaojun Podcast, la transcripción hecha en GPU, la estructura de capítulos y diagramas conceptuales propios\par}
\vspace{0.3cm}
{\large 2025-12-12\par}
\vspace{0.85cm}
\includegraphics[width=0.88\textwidth,height=0.42\textheight,keepaspectratio]{cover.jpg}\par
\vfill
\begin{tcolorbox}[width=0.92\textwidth, colback=black!2!white, colframe=black!55, sharp corners]
\textbf{Título del video}: 123. Entrevista de 3 horas con Wang Guan, fundador de ONE2X: sistemas generativos, sin intermediarios que se queden con el margen y la distribución del poder en las plataformas de contenido\par
\textbf{Canal del video}: Zhang Xiaojun Podcast\par
\textbf{Duración del video}: 03:42:54\par
\textbf{Enlace del video}: \href{\videourl}{\nolinkurl{\videourl}}\par
\textbf{Nota sobre el material}: YouTube no ofrece subtítulos disponibles; el texto se basa en la transcripción de faster-whisper large-v3 en CUDA, los capítulos del video, la portada y figuras didácticas propias. Las tomas fijas de la entrevista no se repiten en el texto.\par
\textbf{Página del proyecto}: \href{\repourl}{\nolinkurl{\repourl}}\par
\end{tcolorbox}
\end{titlepage}

\tableofcontents
\newpage
\section{Introducción: no es una entrevista más sobre una startup de generación de video}

Esta sección explica por qué vale la pena convertir EP123 en un material sobre economía de plataformas y métodos de producto de IA, en lugar de limitarse a registrar la historia emprendedora de ONE2X y Mideo. El hilo conductor de Wang Guan (王冠) tiene, de hecho, tres niveles: primero, cómo cambia el papel del gerente de producto de IA en la era de los modelos grandes; segundo, por qué los datos determinan los límites de la inteligencia y cómo las empresas de aplicaciones encuentran oportunidades mediante los datos que genera el propio producto; tercero, cómo los sistemas generativos reescribirán los sistemas de recomendación, las plataformas de contenido y la distribución del poder entre creadores, plataformas y consumidores.

La frase más contundente del episodio es «sin intermediarios que se queden con el margen». Aquí el intermediario no es una empresa concreta, sino el papel de asignador de inventario que las plataformas de internet han desempeñado durante mucho tiempo: los creadores producen el contenido y la plataforma se encarga de distribuirlo, emparejarlo, ordenarlo, cobrar comisión y monetizarlo. El impacto de los sistemas generativos consiste en que interiorizan la parte de «asignación» dentro de la «producción bajo demanda»: el usuario ya no necesariamente entra primero al acervo de inventario a buscar contenido, sino que entrega su intención al sistema de producción, y este genera directamente contenido que corresponde a lo que el usuario necesita.

\begin{importantbox}{Tesis central del episodio}
En la era de la IA, los bienes de información pasarán de una «economía de plataformas de distribución» a una «economía de sistemas de producción». El contenido dejará de ser solo una obra aislada que entra al fondo de tráfico y podrá convertirse en una IP continua, generada en conjunto por la recipe, el context, el modelo y la intención del usuario; la moneda de las plataformas también podría pasar de la atención a la confianza.
\end{importantbox}

\begin{knowledgebox}{Nota sobre la estrategia visual}
Este video es una entrevista con toma fija, sin slides, pizarrón ni demostración de producto. El texto usa la portada para identificar la fuente, y todas las imágenes del cuerpo son diagramas conceptuales de elaboración propia, que sirven para explicar mecanismos abstractos como las tres etapas de los datos, la pila de un sistema generativo, las diferencias entre sistemas de recomendación y sistemas generativos, y el desplazamiento del poder de las plataformas.
\end{knowledgebox}

\subsection{Resumen de la sección}

El valor de EP123 está en discutir de manera conjunta los productos de IA, el ciclo virtuoso de datos y la estructura de poder de las plataformas de contenido. No se pregunta solo «cómo se hace la generación de video», sino: cuando los sistemas generativos asuman la producción de contenido, ¿cómo cambiarán los papeles de las plataformas, los creadores, los consumidores y los gerentes de producto?
\section{La trayectoria de un gerente de producto: de los datos estructurados a la generación no estructurada}

Esta sección convierte primero la experiencia del entrevistado en pistas metodológicas. La trayectoria profesional de Wang Guan (王冠) no es solo una presentación de antecedentes: explica por qué pasó de preguntarse «cómo ofrecer capacidades algorítmicas» a preguntarse «cómo diseñar las capacidades de un modelo y los sistemas generativos».

Esta sección examina primero la trayectoria profesional de Wang Guan. Abarca big data, CV y aprendizaje profundo, modelos preentrenados, el puesto de gerente de producto de modelos en Moonshot y la fundación de ONE2X. El valor didáctico de esta trayectoria es que ofrece una historia técnica de la IA desde la perspectiva de un gerente de producto: la IA no es solo la evolución de los modelos algorítmicos, sino también la evolución de la forma de representar los datos, del lugar que ocupa el producto y de la manera en que los usuarios perciben la tecnología.

\begin{figure}[H]
\centering
\includegraphics[width=0.96\textwidth]{product-manager-trajectory.png}
\caption{La trayectoria de Wang Guan en productos de IA: big data, CV y aprendizaje profundo, preentrenamiento, Moonshot y ONE2X. Diagrama conceptual propio, elaborado a partir de la conversación entre 00:02:00 y 00:28:39.}
\end{figure}

\begin{knowledgebox}{Lectura de la figura: el hilo común de la trayectoria es «reducir la barrera de entrada para usar capacidades algorítmicas»}
Desde los perfiles de usuario en Baidu, pasando por la plataforma abierta de algoritmos, los frameworks de código abierto y las herramientas de productividad algorítmica, hasta los productos de modelos en Moonshot, el trabajo de Wang Guan siempre ha girado en torno a «permitir que los usuarios accedan a cierta capacidad algorítmica con una barrera de entrada menor». ONE2X es solo la extensión de esa línea hacia la generación de video y los sistemas generativos.
\end{knowledgebox}

\subsection{Las tres eras de producto de la IA}

Esta subsección extrae de esa experiencia un primer juicio: lo decisivo en el cambio de los ciclos de la IA no es solo el tamaño de los modelos, sino qué forma de datos pueden ajustar.

Wang Guan no divide los ciclos de la IA en 1.0, 2.0 y 3.0, sino según si los datos son estructurados o no. El aprendizaje automático tradicional y la CV temprana dependían de etiquetas estructuradas: perfiles de usuario, cuadros delimitadores, etiquetas de categoría y campos explícitos. Los modelos grandes, en cambio, ajustan datos no estructurados: lenguaje, imágenes, video y representaciones de un mundo continuo. Los datos no estructurados tienen mayor capacidad expresiva y se acercan más a la continuidad del mundo real.

\begin{knowledgebox}{Términos clave: datos estructurados y no estructurados}
\begin{tabularx}{\textwidth}{p{0.22\textwidth}p{0.34\textwidth}X}
\toprule
Término & Explicación breve & Relación con este episodio \\
\midrule
Datos estructurados & Datos con campos, etiquetas, coordenadas o categorías explícitos & Entrada central en la era del aprendizaje automático tradicional y de la CV. \\
Datos no estructurados & Representaciones continuas como texto, imágenes y video, que no se comprimen de antemano en campos fijos & Objeto central de entrenamiento de los modelos grandes y de los sistemas generativos. \\
User profile & Perfil de usuario: describe las preferencias y el comportamiento del usuario mediante etiquetas & Base del modelado del lado del usuario en los sistemas de recomendación. \\
Gerente de producto de modelos & Diseña el desempeño, las capacidades y la evaluación del modelo para que el usuario perciba esas capacidades & Puesto clave que se formó durante la experiencia en Moonshot. \\
\bottomrule
\end{tabularx}
\end{knowledgebox}

\subsection{Por qué los modelos no han reemplazado al gerente de producto}

Wang Guan sostiene que en la era de la IA el valor del gerente de producto ha aumentado, no disminuido. Hay dos razones. Primero, el modelo se parece al System1 humano: comprime de antemano una gran cantidad de información dentro del sistema y puede reaccionar a la entrada de forma rápida e instintiva. El gerente de producto debe diseñar qué capacidades debe tener ese System1, y lograr que esas capacidades se entrenen y que el usuario las perciba. Segundo, la salida del modelo depende de qué tokens la preceden; el gerente de producto debe diseñar el workflow, el framework del agent, la base de conocimiento del dominio y el context, para que el modelo obtenga más tokens útiles al generar.

\begin{importantbox}{Las dos nuevas tareas del gerente de producto de IA}
Primero, definir las capacidades del modelo: qué desempeño vale la pena entrenar, cómo evaluarlo y cómo lo percibe el usuario. Segundo, diseñar el context: cómo convertir el know-how del negocio, los flujos de trabajo, las herramientas y el conocimiento del dominio en tokens útiles que el modelo pueda aprovechar.
\end{importantbox}

\subsection{Resumen de la sección}

La trayectoria de Wang Guan muestra que el gerente de producto de IA no queda absorbido por las capacidades del modelo, sino que pasa a un primer plano. Cuanto más capaz es el modelo, más se necesita a alguien que defina sus capacidades, organice el context, diseñe la forma de las tareas y convierta esas capacidades en un producto que los usuarios puedan entender, por el que estén dispuestos a pagar y que puedan usar de forma continua.
\section{La compresión es inteligencia: por qué el lenguaje se volvió el soporte de la inteligencia}

La sección anterior trató cómo entiende el modelo un gerente de producto; esta sección aborda uno de los primeros principios de este episodio: la compresión es inteligencia. Según Wang Guan (王冠), la compresión une datos que antes eran puntos discretos en una estructura continua, y esa continuidad se manifiesta hacia fuera como generalización, emergencia, alucinación o inteligencia. El lenguaje importa porque puede expresar un mundo suficientemente rico con un costo de entrenamiento relativamente bajo.

\begin{figure}[H]
\centering
\includegraphics[width=0.96\textwidth]{compression-intelligence.png}
\caption{La compresión es inteligencia: la compresión une tareas discretas en una capacidad generalizable. Diagrama conceptual propio, elaborado a partir de la conversación de 00:28:39--00:32:25.}
\end{figure}

\begin{knowledgebox}{Lectura de la figura: entre los datos y la inteligencia no hay magia, sino continuidad}
En la figura, los «datos» son al principio muestras discretas; la compresión hace que formen relaciones continuas, el lenguaje aporta un soporte de alta capacidad expresiva, la generalización se manifiesta como la capacidad de resolver tareas no vistas y, al final, las personas llaman a eso inteligencia. Esta cadena explica por qué un modelo puede combinar «traducción del chino» y «resumen en chino» para obtener una capacidad como «resumen en inglés», que no se entrenó de forma explícita.
\end{knowledgebox}

\subsection{Por qué la compresión produce generalización}

En el programa se da un ejemplo intuitivo: si el modelo se entrenó en traducción del chino al inglés y también en resumen en chino, pero no en resumen en inglés, la representación comprimida puede permitirle unir las tareas de «traducir» y «resumir» y, así, producir un resumen en inglés. Esta capacidad no surge de la nada: en la distribución de los datos, las distintas tareas forman una estructura continua y transferible.

\begin{warningbox}{No hay que malinterpretar «la compresión es inteligencia» como una explicación universal}
La compresión explica una parte del origen de la generalización, pero eso no significa que cualquier compresión produzca inteligencia útil. La calidad de los datos, el objetivo de entrenamiento, la arquitectura del modelo, la capacidad de cómputo, la evaluación y el context del producto influyen en la capacidad final.
\end{warningbox}

\subsection{Implicaciones de producto de «el lenguaje es inteligencia»}

Más adelante, Wang Guan actualizó «la compresión es inteligencia» a «el lenguaje es inteligencia». Esta afirmación no niega las imágenes, el video ni las acciones; subraya que el lenguaje, como soporte abstracto, puede expresar a bajo costo el mundo, las tareas, los métodos, las relaciones y las intenciones. Para el producto, esto significa que el lenguaje natural no es solo una entrada de la UI, sino también un DSL, una descripción de tareas, un criterio de evaluación y una forma de organizar el context.

\subsection{Resumen de la sección}

La idea de que la compresión es inteligencia es la base de todos los juicios posteriores de este episodio: los datos determinan los límites, el lenguaje comprime el mundo en expresiones combinables, y el gerente de producto debe diseñar tanto las capacidades que el modelo puede aprender como el context que el modelo puede aprovechar.
\section{Los datos como límite de la inteligencia: datos públicos, datos de dominio y datos generados por el producto}

Las dos secciones anteriores explican cómo se obtienen las capacidades por compresión; esta sección analiza qué determina el límite de esas capacidades. Wang Guan (王冠) reformula «tanta inteligencia como trabajo humano haya» como «tanta inteligencia como datos haya». Aquí los datos no son solo archivos o tablas: son la comprensión del problema, el know-how del negocio, el proceso de las tareas y la retroalimentación de los usuarios.

\begin{figure}[H]
\centering
\includegraphics[width=0.96\textwidth]{data-three-stages.png}
\caption{Las tres etapas de los datos: los datos públicos, los datos de dominio y los datos generados por el producto corresponden a distintos actores. Diagrama conceptual de elaboración propia, basado en la conversación de 00:37:11--01:05:36.}
\end{figure}

\begin{knowledgebox}{Lectura de la figura: cada etapa favorece a actores distintos}
La etapa de los datos públicos favorece a las empresas de modelos base, porque todas limpian y comprimen datos de internet; la etapa de los datos de dominio favorece a las grandes empresas tecnológicas y a las industrias tradicionales, porque poseen escenarios de negocio y datos privados; solo la etapa de los datos generados por el producto se parece más a una oportunidad para emprender con aplicaciones, porque un producto nuevo puede generar, al usarse, datos que antes no existían.
\end{knowledgebox}

\subsection{Una visión espacio-temporal de la inteligencia: datos, cómputo y algoritmos}

Wang Guan explica el límite de la inteligencia con una analogía en un espacio bidimensional. Los datos determinan la frontera del círculo, es decir, el alcance de capacidades al que el modelo puede llegar finalmente; el cómputo determina la velocidad con que se aproxima a esa frontera; los algoritmos determinan la trayectoria y la eficiencia de esa aproximación. Esta postura no niega la importancia de los algoritmos, sino que los sitúa dentro del límite que fijan los datos: sin los datos correspondientes, por potente que sea el algoritmo, difícilmente se obtienen capacidades estables de la nada.

\begin{importantbox}{Los datos determinan el límite de la inteligencia}
Si una tarea requiere know-how del sector, escenarios reales y juicio ante problemas abiertos, entonces «conseguir más datos» normalmente no es trabajo menor, sino que equivale a definir el problema mismo. Los gerentes de producto son importantes porque pueden convertir la comprensión del negocio en datos de los que se puede aprender.
\end{importantbox}

\subsection{Por qué las dos primeras etapas no son oportunidades típicas para emprender con aplicaciones}

Esta sección relaciona además las tres etapas de los datos con las oportunidades de emprendimiento. No todo tipo de datos está al alcance de una empresa emergente, ni toda etapa es adecuada para que entre una empresa de aplicaciones.

En la etapa de los datos públicos se compite en densidad de talento, cómputo, eficiencia de ingeniería y desgaste organizacional, por lo que resulta más adecuada para las empresas de modelos base. La etapa de los datos de dominio favorece a las grandes empresas tecnológicas y a las industrias con un alto grado de digitalización, porque ya cuentan con escenarios de negocio, canales de clientes y datos privados. La verdadera oportunidad para las empresas emergentes de aplicaciones está en la tercera etapa: los datos generados por el producto, es decir, los datos que surgen mientras el producto funciona, que antes no existían y que no pueden comprarse directamente fuera.

\begin{knowledgebox}{Ejemplos de datos generados por el producto}
\begin{tabularx}{\textwidth}{p{0.22\textwidth}p{0.34\textwidth}X}
\toprule
Tipo de datos & Cómo se generan & Por qué constituyen una barrera de entrada \\
\midrule
Trayectorias de intención del usuario & El usuario describe, corrige y elige continuamente dentro del producto & Reflejan necesidades y preferencias reales. \\
Muestras de generación fallida & El proceso de modificación cuando el contenido generado no satisface & Exponen las debilidades del modelo y los escenarios del producto. \\
Recipe & Métodos de producción y combinaciones de parámetros reutilizables & No son contenido aislado, sino capacidad de producción. \\
Retroalimentación del ciclo comercial & Si el contenido se reproduce, se paga, convierte o se reutiliza & Se vincula directamente con el valor del producto. \\
\bottomrule
\end{tabularx}
\end{knowledgebox}

\subsection{Resumen de la sección}

Las tres etapas de los datos explican por qué las empresas de aplicaciones no pueden limitarse a una capa delgada de UI. La verdadera oportunidad para las aplicaciones está en crear datos nuevos: cuanto más se usa el producto, más intenciones, procesos, fallas, recipes y retroalimentación comercial genera que no existen fuera de él.
\section{Por qué elegir la generación de video: una modalidad de alto valor en una ventana de inmadurez}

La sección anterior explicó que las oportunidades de aplicación surgen de los datos endógenos del producto; esta sección examina por qué ONE2X eligió la generación de video. Wang Guan (王冠) eligió el video no porque fuera lo más fácil, sino porque tiene un valor alto, la modalidad aún no ha madurado y su espacio de comercialización es evidente; además, el procesamiento de video no se limita a generar un solo video, sino que puede convertirse en un sistema completo articulado en torno al DSL, el context y el flujo de producción.

\begin{figure}[H]
\centering
\includegraphics[width=0.96\textwidth]{video-generation-bet.png}
\caption{Por qué elegir la generación de video: modalidad de alto valor, ventana de inmadurez, espacio de producto, validación de ingresos y reordenamiento de las plataformas. Diagrama conceptual propio, elaborado a partir de la conversación de 01:05:36--01:26:15.}
\end{figure}

\begin{knowledgebox}{Lectura de la figura: la generación de video no consiste en «hacer un botón de generar»}
En la figura aparecen a la vez el alto valor, la inmadurez, el espacio de procesamiento, la validación de ingresos y el reordenamiento de las plataformas. Lo que Wang Guan valora es el alto valor del video como bien de información, así como el impacto a largo plazo de los sistemas de generación sobre la producción de contenido, los flujos de trabajo de los creadores y el poder de distribución de las plataformas.
\end{knowledgebox}

\subsection{Por qué la modalidad de video tiene valor comercial}

Esta subsección explica el video primero desde su valor comercial, no desde el lucimiento técnico. Para un equipo en etapa temprana, la elección de modalidad debe responder primero: por qué el usuario estaría dispuesto a pagar y en qué escenario puede validarse el valor lo antes posible.

El video es más pesado que el texto y la imagen: su costo de producción es mayor, su ancho de banda de consumo es más amplio y está más cerca de la distribución en plataformas y de la conversión comercial. Incluso cuando la tecnología de generación aún no ha madurado, las herramientas de procesamiento de video y el SaaS ya pueden generar ingresos considerables. Para un equipo en etapa temprana, elegir una modalidad de alto valor significa validar antes la disposición a pagar y el valor para los creadores.

\subsection{El criterio de producto de Mideo}

A continuación se pasa de la elección de modalidad a la forma del producto. Lo esencial de Mideo no es «si puede generar un fragmento de video», sino si puede consolidar el método de producción de cierto tipo de contenido.

En el programa se insiste repetidamente en que lo que hace ONE2X no es simplemente generate video, sino un sistema de procesamiento y producción de video en un sentido más amplio. El producto debe comprender el production method de cierta categoría, abstraerlo en un DSL y una recipe, y después hacer que el modelo genere contenido continuo que cumpla el objetivo. Wang Guan mencionó que el equipo incluso ganó dinero primero con su propio producto, lo que demostró que puede sostener la producción de los creadores y la monetización en las plataformas.

\begin{importantbox}{De la capacidad de generación al sistema de producción}
Generar un video es solo una capacidad del modelo; producir de forma estable cierto tipo de video, reutilizar el método, iterar la recipe y cerrar el ciclo comercial es lo que constituye un sistema de producción.
\end{importantbox}

\subsection{Resumen de la sección}

La oportunidad de la generación de video proviene de una modalidad de alto valor y de una ventana de inmadurez, pero la verdadera barrera competitiva no está en «poder generar», sino en lograr abstraer el método de producción de video en un sistema reutilizable y en generar, con el uso, datos endógenos del producto.
\section{Pila del sistema de generación: DSL, Context, Agent/Workflow y System1}

Después de discutir por qué hacer video, esta sección entra en la metodología central que propone Wang Guan (王冠): el sistema de generación. Un sistema de generación no es un solo modelo ni un solo workflow, sino un método de producción organizado en torno a cierto tipo de bien de información. Integra en un sistema la expresión del problema, el context efectivo, el agent/workflow, las capacidades del modelo y la evaluación.

\begin{figure}[H]
\centering
\includegraphics[width=0.92\textwidth]{generation-system-stack.png}
\caption{Pila del sistema de generación: System1, DSL, Context, Agent/Workflow y evaluación trabajan en conjunto. Diagrama conceptual propio, elaborado a partir de la conversación de 02:01:44--02:25:49.}
\end{figure}

\begin{knowledgebox}{Lectura de la figura: el sistema de generación es principalmente System2}
System1 es el modelo base; el DSL sirve para expresar el problema y el método de producción; el Context aporta tokens efectivos; el Agent/Workflow organiza los pasos; la evaluación mide el grado de inteligencia del sistema. Wang Guan considera que agent y workflow no tienen por qué oponerse: en esencia, ambos producen un context más efectivo para que System1 genere mejores salidas.
\end{knowledgebox}

\subsection{DSL: expresar el método de producción}

DSL es Domain-Specific Language, lenguaje específico de dominio. Aquí no se refiere necesariamente a un lenguaje de programación tradicional, sino al sistema de expresión de una tarea de generación. En el caso del video, el DSL debe poder describir tomas, ritmo, imagen, estilo, material, transiciones, público objetivo y método de producción. Sin DSL, el sistema solo puede acumular prompts; con DSL, el sistema puede reutilizar el método de producción de forma estable.

\subsection{Context Engineering: el nuevo campo principal del gerente de producto}

Esta sección pasa del DSL al context, porque el resultado de un sistema de generación no depende solo del modelo, sino también de lo que el modelo ve antes de generar y de cómo entiende la estructura de la tarea.

Wang Guan considera muchos agents, workflows, bases de conocimiento y repositorios especializados como mecanismos para generar context. La salida del modelo depende de los tokens previos; lo que el producto debe hacer es convertir el know-how del negocio en un context que el modelo pueda entender, que sea económicamente rentable y que dé resultados estables. Este sistema no es necesariamente el más complejo en lo técnico, pero depende en gran medida de la comprensión del negocio y del criterio de producto.

\begin{warningbox}{No subestimar el Workflow como algo «sin contenido técnico»}
Workflow, Agent, repositorios especializados y DSL sirven para proporcionar al modelo un context efectivo. Su valor no reside solo en la complejidad de ingeniería, sino en si logran expresar de forma estable los métodos del negocio, la estructura de la tarea y los objetivos de evaluación.
\end{warningbox}

\subsection{Métrica de estrella polar: el grado de inteligencia del sistema}

Esta sección responde «qué optimiza en realidad un producto de IA». El valor comercial es el mínimo indispensable, pero si solo se mira el ingreso a corto plazo, se pierde de vista el crecimiento del sistema de generación como capacidad de producción en sí misma.

Wang Guan sostiene que un producto AI Native debe tener valor comercial, desde luego, pero que la métrica de estrella polar de nivel superior es el «grado de inteligencia del sistema». No se trata de algo esotérico, sino de preguntar: ¿el sistema puede producir de forma automática, elegir métodos de forma automática, corregirse de forma automática y llegar por sí mismo a mejores salidas? En un sistema de generación, el usuario no solo compra una herramienta, sino que usa un sistema de producción que cada vez sabe hacer mejor su trabajo.

\begin{figure}[H]
\centering
\includegraphics[width=0.96\textwidth]{ai-product-north-star.png}
\caption{Estrella polar de un producto de IA: por encima del valor comercial, medir el producto por el grado de inteligencia del sistema. Diagrama conceptual propio, elaborado a partir de la conversación de 02:50:34--03:05:45.}
\end{figure}

\begin{knowledgebox}{Lectura de la figura: el grado de inteligencia no es un lema abstracto}
La figura va del valor comercial al diseño de capacidades, y de ahí al context, la producción automática y el grado de inteligencia. El grado de inteligencia del sistema puede manifestarse en: si entiende el objetivo, si elige el método adecuado, si reduce los pasos manuales, si puede mejorar a partir de sus fallas y si permite al usuario lograr una producción de mayor calidad a menor costo.
\end{knowledgebox}

\subsection{Resumen de la sección}

El sistema de generación es el núcleo con el que una empresa de aplicaciones puede construir una ventaja defensiva. Convierte las capacidades del modelo en métodos de producción, el know-how del negocio en context y cada generación aislada en una recipe reutilizable, y mide el avance del producto por el grado de inteligencia del sistema.
\section{La frontera entre las empresas de aplicaciones y las empresas de modelos se volverá difusa}

La sección anterior trató la pila del sistema de generación; esta sección analiza las fronteras entre empresas. Wang Guan (王冠) considera que la frontera entre las empresas de aplicaciones y las empresas de modelos se volverá difusa. La razón es que las empresas de modelos subirán hacia el producto, mientras que las empresas de aplicaciones bajarán hacia los datos, el DSL, el context y la optimización de modelos locales. La diferencia real no está en la etiqueta de la empresa, sino en quién posee el volante de datos y los escenarios de usuario más decisivos.

\begin{figure}[H]
\centering
\includegraphics[width=0.94\textwidth]{app-model-boundary.png}
\caption{Frontera entre las empresas de aplicaciones y las empresas de modelos: el modelo base y los datos generados por el propio producto forman volantes distintos. Diagrama conceptual propio, elaborado a partir de la conversación de 01:41:59--02:01:44.}
\end{figure}

\begin{knowledgebox}{Lectura de la figura: una frontera difusa no es una frontera que desaparece}
Las empresas de modelos poseen el modelo base, la capacidad de cómputo, las capacidades generales y la distribución por API; las empresas de aplicaciones poseen los escenarios, los datos generados por el propio producto, los flujos de trabajo y el valor para el usuario. Lo decisivo en la competencia futura será quién logre conectar las capacidades del modelo con los datos del escenario en un ciclo cerrado.
\end{knowledgebox}

\subsection{Por qué las empresas de modelos suben hacia el producto}

Esta subsección examina primero por qué las empresas de modelos se desplazan hacia la capa de aplicaciones. Si las capacidades del modelo se quedan solo en la API, quedan demasiado lejos de la retroalimentación de los usuarios y del techo de ingresos; un punto de entrada de producto puede convertir las capacidades del modelo en una relación comercial de mayor valor.

Si una empresa de modelos solo vende API, queda fácilmente restringida por el precio y el costo de cómputo; subir hacia el producto le permite obtener el punto de entrada de los usuarios, retroalimentación real e ingresos más altos. ChatGPT, Claude Code, Gemini y otros demuestran que las empresas de modelos entran directamente en la capa de aplicaciones. Sobre todo en escenarios como Coding, las empresas de modelos ya controlan de origen los datos y el punto de entrada de los desarrolladores, por lo que les resulta más fácil hacerlo por su cuenta.

\subsection{Por qué las empresas de aplicaciones deben bajar hacia los modelos y los datos}

A continuación se examina el mismo problema desde la perspectiva de las empresas de aplicaciones. Si las empresas de modelos suben hacia el producto, las empresas de aplicaciones no pueden quedarse en la capa superficial, sino que deben construir hacia abajo datos, métodos y flujos de trabajo.

Si una empresa de aplicaciones solo hace la UI, también será copiada por las empresas de modelos. Debe construir hacia abajo el DSL, el context, los datos del dominio, la evaluación, la retroalimentación y la optimización local. No necesariamente tiene que entrenar un gran modelo general, pero sí necesita que su sistema de generación cuente con datos generados por el propio producto que otros no puedan obtener.

\begin{importantbox}{La ventaja defensiva de las empresas de aplicaciones}
La ventaja defensiva de una empresa de aplicaciones no es solo la interfaz, sino «el volante de datos dentro del escenario»: la intención del usuario, el proceso de producción, los ejemplos fallidos, las recipe, la retroalimentación comercial y la evaluación continua.
\end{importantbox}

\subsection{Resumen de la sección}

La frontera entre las empresas de aplicaciones y las empresas de modelos se volverá difusa, pero las diferencias no desaparecerán. Las empresas de modelos parten de las capacidades; las empresas de aplicaciones, de los escenarios. Quien logre formar un volante de datos más sólido obtendrá una posición más estable en la era de los sistemas de generación.
\section{Sistemas de recomendación vs.\ sistemas generativos: sin intermediarios que se queden con la diferencia}

Esta sección entra en la teoría de plataformas más central de todo el episodio. La esencia de un sistema de recomendación es la asignación de inventario: la plataforma empareja el contenido existente con las preferencias de los usuarios y controla el ordenamiento, la distribución y la fijación del precio del tráfico. La esencia de un sistema generativo es la producción bajo demanda: la necesidad del usuario entra directamente en el extremo de producción, que genera el contenido y lo entrega directamente al usuario; la etapa de asignación queda interiorizada en el sistema de producción.

\begin{figure}[H]
\centering
\includegraphics[width=0.94\textwidth]{recommendation-vs-generation.png}
\caption{Sistemas de recomendación vs.\ sistemas generativos: la diferencia de fondo entre la asignación de inventario y la producción bajo demanda. Diagrama conceptual propio, elaborado a partir de la conversación de 02:38:11--02:50:34.}
\end{figure}

\begin{knowledgebox}{Lectura de la figura: el intermediario es el «derecho de asignación del inventario»}
El sistema de recomendación empareja contenido a partir del inventario existente, y la plataforma controla el derecho de distribución; el sistema generativo produce al instante según la necesidad del usuario, y la asignación queda interiorizada en la producción. La expresión «sin intermediarios que se queden con la diferencia» no significa que no haya plataforma, sino que el derecho de distribución de las plataformas tradicionales deja de ser la única forma de convertir el valor en ingresos.
\end{knowledgebox}

\subsection{Por qué el sistema de recomendación es un intermediario}

Este apartado explica primero el sentido mecánico de la palabra «intermediario». No es una crítica moral: significa que la plataforma controla la asignación del inventario, el ordenamiento y el poder de comercialización, y por eso puede extraer valor entre productores y consumidores.

El poder de las plataformas de internet proviene de una relación bilateral: los creadores producen contenido, los consumidores lo consumen y la plataforma se encarga del emparejamiento y extrae valor. Los motores de búsqueda, los motores de recomendación y las plataformas de comercio electrónico actúan, en distinto grado, como intermediarios. Los creadores invierten recursos en producir contenido, pero el mismo contenido publicado desde cuentas distintas puede recibir un tráfico completamente diferente; esto muestra que el mecanismo de asignación no depende solo del contenido en sí, sino también del peso que la plataforma asigna, del historial de la cuenta, de los objetivos del algoritmo y de las reglas comerciales.

\begin{warningbox}{No hay que entender «la plataforma es un intermediario» como que la plataforma no aporta valor}
Las plataformas sí crean valor en descubrimiento, emparejamiento, prevención de fraude, pagos, comunidad y distribución a escala. El problema es que los sistemas generativos reducen la escasez relativa de la «asignación de inventario» y obligan a las plataformas a desplazarse hacia el sistema de producción, las redes de confianza u otras capas de valor.
\end{warningbox}

\subsection{Cómo interioriza la asignación el sistema generativo}

El sistema generativo no busca en el inventario el contenido más cercano a la necesidad del usuario, sino que produce el contenido directamente a partir de la intención del usuario. Sigue habiendo asignación, pero ocurre dentro del proceso de generación: la necesidad y las preferencias del usuario, el context, la recipe y las capacidades del sistema determinan en conjunto la salida. Así, la plataforma deja de ser la única que controla si un contenido se ve o no, y la intención del usuario se convierte en una entrada más directa de la producción.

\begin{figure}[H]
\centering
\includegraphics[width=0.96\textwidth]{power-shift-to-consumer.png}
\caption{El poder se desplaza hacia el lado del consumidor: la cadena de los bienes de información pasa de la asignación por la plataforma a estar impulsada por la intención del usuario. Diagrama conceptual propio, elaborado a partir de la conversación de 02:25:49--02:38:11.}
\end{figure}

\begin{knowledgebox}{Lectura de la figura: el consumidor no es solo el extremo receptor}
En la cadena tradicional, el productor produce, la plataforma distribuye y el consumidor recibe. El sistema generativo hace que la intención del consumidor entre en el extremo de producción: el usuario no solo consume contenido de forma pasiva, sino que, al consumirlo, participa en su producción.
\end{knowledgebox}

\subsection{Resumen de la sección}

La diferencia de fondo entre los sistemas de recomendación y los sistemas generativos es la diferencia entre la asignación de inventario y la producción bajo demanda. Las plataformas no desaparecerán, pero la forma en que las plataformas tradicionales extraen valor mediante el derecho de distribución se verá cuestionada.
\section{De la plataforma de distribución a la plataforma de producción y distribución: Recipe, IP y confianza}

La sección anterior trató la recomendación y la generación; esta sección sigue proyectando la forma de las plataformas. Wang Guan (王冠) considera que las plataformas futuras pasarán de ser plataformas de distribución a ser plataformas de producción y distribución. Antes se producía contenido puntual, que entraba en un fondo común de tráfico para que el algoritmo lo emparejara con su público. En el futuro, los productores podrían producir recipes, es decir, métodos de producción reutilizables, y a partir de ellos formar contenido continuo e IP. La moneda de la economía de plataformas también pasará de la atención a la confianza.

\begin{figure}[H]
\centering
\includegraphics[width=0.96\textwidth]{distribution-to-production.png}
\caption{De la plataforma de distribución a la plataforma de producción y distribución: contenido puntual, Recipe, IP y confianza como moneda. Diagrama conceptual propio, elaborado a partir de la conversación de 03:20:18--03:42:54.}
\end{figure}

\begin{knowledgebox}{Lectura de la figura: la escasez no desaparece, solo se desplaza}
En la era de las plataformas de distribución, la atención es la moneda y el contenido puntual entra en el fondo común de tráfico a competir por exposición. En la era de las plataformas de producción y distribución, la recipe permite producir contenido de forma continua, la IP sostiene un estilo y una calidad estables, y la confianza se convierte en una moneda de mayor valor.
\end{knowledgebox}

\subsection{Recipe: del contenido puntual al método de producción}

Esta sección traslada el cambio en la forma de las plataformas al objeto que producen los creadores. Una vez que la generación con IA reduce el costo de producir contenido puntual, lo verdaderamente valioso pasará de «esta pieza de contenido» a «el método para producir de forma sostenida este tipo de contenido».

Una recipe no es el video A ni el video B, sino el método de producción común que hay detrás de ambos. Puede incluir la elección de temas, la estructura, las tomas, el ritmo, el estilo, los materiales, el público objetivo y la forma de interacción. Quien produce una recipe no solo produce piezas de contenido aisladas, sino una capacidad reutilizable de producción de contenido.

\begin{importantbox}{La importancia de la recipe}
Cuando el contenido puede generarse a bajo costo, la escasez del contenido puntual disminuye; la escasez se desplaza hacia el método de producción, el gusto, la continuidad y la confianza. La recipe es la manera de estructurar esa escasez.
\end{importantbox}

\subsection{De la atención a la confianza}

Esta sección analiza el cambio de «moneda» en la economía de plataformas. La atención sigue siendo importante, pero cuando la oferta de contenido se expande sin límite, la razón por la que un usuario sigue a un autor o a una IP durante mucho tiempo vale más que un solo clic.

La moneda de las plataformas tradicionales es la atención: reproducciones, tiempo de permanencia, clics e interacciones. Wang Guan considera que, cuando la cantidad de contenido crece enormemente, la atención no alcanza para repartirse y el valor se desplaza hacia la confianza. Substack, Medium, OnlyFans, Zhishi Xingqiu (知识星球) y las comunidades privadas, entre otros, ya muestran esta tendencia: los usuarios no consumen porque la plataforma les recomiende algo, sino que pagan porque confían en un autor, un canal, una IP o un método de producción.

\begin{warningbox}{La confianza no sustituye a la atención: es una moneda de nivel superior}
La atención sigue existiendo, pero la confianza puede sostener un mayor gasto por cliente, relaciones de más largo plazo y públicos más precisos. Cuando la generación con IA aumenta la oferta de contenido, el gusto confiable y la IP continua se vuelven, en realidad, más escasos.
\end{warningbox}

\subsection{¿Desaparecerán los creadores?}

Wang Guan no cree que los creadores vayan a desaparecer; cree que el grupo de creadores se ampliará y que la forma de su valor cambiará. El propio consumo podría convertirse en creación: cuando los usuarios usan GPT-5 o un sistema de generación para resolver problemas, producen contenido valioso para otras personas. El creador no necesariamente edita cada cuadro ni escribe cada frase en persona, sino que participa en la producción mediante la intención, el gusto, la recipe y la confianza.

\subsection{Resumen de la sección}

Que las plataformas pasen de la distribución a la producción y distribución significa que el activo central de la economía del contenido se desplaza del contenido puntual y la distribución por parte de la plataforma hacia la recipe, la IP, la confianza y la intención del usuario. Los creadores no desaparecerán sin más: pasarán de ser productores de contenido a ser organizadores de métodos de producción, de gusto y de confianza.
\section{Organización: ingenieros integrales de IA, colaboración remota y alta densidad de talento}

Después de tratar el producto y la plataforma, esta sección examina la forma de organización de ONE2X. Wang Guan (王冠) concibe ONE2X como un estudio de productos de la era de la IA, no como una empresa tradicional. El estudio pone el énfasis en la exploración, el interés, la calidad y los superindividuos. Cada integrante del equipo debe poder hacerse cargo de un área por sí solo; a estas personas se les llama ingenieros integrales de IA.

\subsection{Por qué un equipo pequeño puede hacer más}

Esta subsección pasa de la teoría de la plataforma a la ejecución organizacional. Un sistema generativo no se sostiene solo con grandes narrativas: necesita un equipo pequeño capaz de validar, producir y aprender con rapidez para llevar el método a la práctica.

En la era de la IA, buena parte de la capacidad productiva existe en forma de token, lo que equivale a que el equipo puede recurrir a mano de obra externa. Un equipo pequeño con alta densidad de talento, si tiene una fuerte motivación propia y una dirección clara, puede lograr resultados que antes exigían un equipo mucho mayor. En la contratación, ONE2X se fija sobre todo en la motivación propia, la iniciativa, la construcción en público, los proyectos de código abierto, la difusión de lo aprendido y en si la reward function personal tiene proyección sobre la dirección del equipo.

\begin{knowledgebox}{Palabras clave de la organización}
\begin{tabularx}{\textwidth}{p{0.22\textwidth}p{0.34\textwidth}X}
\toprule
Palabra clave & Significado & Función en la organización \\
\midrule
Ingeniero integral de IA & Persona capaz de entender los problemas desde el producto, la ingeniería, los modelos, el contenido y el negocio & Adecuado para un equipo pequeño en etapa de exploración. \\
Superindividuo & Persona muy motivada por sí misma, capaz de producir de forma independiente y de amplificar sus capacidades con IA & Reduce los costos de gestión y de coordinación. \\
Reward function & Lo que realmente motiva a la persona & Al contratar, debe apuntar en el mismo sentido que los objetivos del equipo. \\
Remote & Trabajo a distancia & Amplía el conjunto de talento disponible, pero obliga a resolver la soledad y la confianza. \\
\bottomrule
\end{tabularx}
\end{knowledgebox}

\subsection{Problemas y soluciones del trabajo a distancia}

El trabajo a distancia trae consigo sensación de soledad y problemas de confianza, porque los mensajes escritos carecen de tono, de lenguaje corporal y de contexto. Wang Guan no renunció por ello al modelo remote; lo trata como una restricción organizacional: contrarresta esos problemas con documentación, comunicación asíncrona, una fuerte motivación propia y objetivos claros. El trabajo presencial también tiene sus propios problemas; el trabajo a distancia no es más cómodo, sino otra manera de obtener un conjunto de talento más amplio y mayor autonomía.

\subsection{Resumen de la sección}

La forma organizacional de ONE2X está al servicio de la exploración de sistemas generativos: equipo pequeño, alta densidad de talento, fuerte motivación propia, colaboración remota y productividad amplificada por la IA. No es una respuesta universal, pero es coherente con su definición como «estudio de productos».
\section{Términos clave: índice de conceptos de este episodio}

Este episodio maneja muchos conceptos, y varios de estos términos no se usan en el sentido habitual de los libros de texto. Aquí se reúnen para facilitar el repaso.

\begin{longtable}{p{0.24\textwidth}p{0.34\textwidth}p{0.33\textwidth}}
\toprule
Término & Explicación en una frase & Papel en este episodio \\
\midrule
ONE2X & Estudio de productos de IA fundado por Wang Guan (王冠) & Caso práctico de este episodio. \\
Mideo & Producto de ONE2X para generar y procesar video con IA & Sirve para validar la generación de video y el sistema de generación. \\
System1 & Capacidad de respuesta rápida e instintiva que el modelo comprime & El gerente de producto debe diseñar las capacidades del modelo y su evaluación. \\
System2 & Razonamiento, flujos de trabajo, context y sistema de generación organizados en torno al modelo & Terreno principal donde las empresas de aplicaciones construyen sus barreras competitivas. \\
Context Engineering & Diseño del contexto útil que el modelo tiene disponible antes de generar & Convierte el know-how del negocio en tokens que el modelo puede usar. \\
DSL & Lenguaje o sistema de expresión específico de un dominio & Sirve para describir los métodos de producción de video y la estructura de los problemas. \\
Recipe & Método de producción reutilizable, no una pieza de contenido aislada & Elemento decisivo para pasar de la producción de contenido a la IP y la confianza. \\
Datos generados por el producto & Datos que surgen del uso del producto y que antes no existían & Ahí está la oportunidad para las empresas emergentes de aplicaciones. \\
Sistema de generación & Sistema que organiza la producción de contenido según la intención y la recipe & Desafía el poder de plataforma de los sistemas de recomendación. \\
Sistema de recomendación & Asigna a cada usuario contenido tomado de un inventario existente & Base del poder de distribución de las plataformas de internet. \\
Plataforma de distribución & Plataforma que controla la asignación de contenido y la distribución del tráfico & Forma de las plataformas de la internet anterior. \\
Plataforma de producción y consumo & Plataforma que une producción y consumo y genera contenido según la intención del usuario & Nueva forma que podría surgir en la era de la IA. \\
Moneda de atención & La plataforma pone precio a la atención: reproducciones, tiempo de permanencia, clics & Moneda central de las plataformas de contenido tradicionales. \\
Moneda de confianza & El usuario paga o consume de forma continua porque confía en un autor, una IP o un canal & Valor que podría volverse más escaso cuando la IA genere el contenido. \\
AGI en sentido amplio & Narrativa de largo plazo que, a escala macro, conduce a la inteligencia general & En el programa se usa para distinguirla de las oportunidades concretas de producto. \\
AGI en sentido estricto & Sistema inteligente suficientemente general dentro de un tipo de tarea o escenario & Más cercana a una ruta que las empresas de aplicaciones pueden validar. \\
\bottomrule
\end{longtable}

\subsection{Resumen de la sección}

En conjunto, estos términos forman un marco de economía de plataformas: los datos determinan el límite de la inteligencia, el context libera las capacidades del modelo, el sistema de generación reescribe la producción de contenido, y la recipe y la confianza sustituyen a las piezas de contenido aisladas y a la atención como los activos más importantes.
\section{Síntesis y ampliación}

\subsection{Conclusiones centrales}

Esta sección condensa todo el texto en criterios reutilizables. El lector puede usar estas conclusiones como lista de verificación para determinar si una aplicación de IA tiene realmente un sistema generativo, datos generados dentro del propio producto y potencial para reestructurar plataformas.

Para cerrar, EP123 se condensa en diez conclusiones.

\begin{enumerate}
\item Los modelos no han sustituido al gerente de producto de IA; este pasa al primer plano en la definición de las capacidades del modelo y el diseño del context.
\item La compresión hace que tareas discretas formen una estructura continua; el lenguaje es un portador importante de inteligencia que expresa el mundo a bajo costo.
\item Los datos determinan los límites de la inteligencia, la capacidad de cómputo determina la velocidad con que uno se acerca a ellos y los algoritmos determinan la eficiencia del camino.
\item Los datos de dominio público favorecen a los modelos base, los datos de dominio específico favorecen a las grandes empresas tecnológicas y los datos generados dentro del producto favorecen a los emprendimientos de aplicaciones.
\item La oportunidad en la generación de video no está en un solo botón de generate, sino en el método de producción, la recipe y el ciclo comercial cerrado.
\item Un sistema generativo se compone de DSL, context, agent/workflow, modelos y evaluación en conjunto.
\item La frontera entre empresas de aplicaciones y empresas de modelos se volverá difusa, pero sus ciclos virtuosos de datos son distintos.
\item Un sistema de recomendación asigna inventario; un sistema generativo produce bajo demanda; el poder de distribución de las plataformas tradicionales se debilitará.
\item La economía del contenido pasará del contenido individual y la atención a la recipe, la IP y la confianza.
\item Los creadores no desaparecerán sin más: pasarán de ser productores de contenido a ser organizadores de métodos de producción, gusto y confianza.
\end{enumerate}

\subsection{Preguntas abiertas}

Al final se dejan preguntas abiertas, porque muchos juicios de este episodio aún están en fase de proyección. Los sistemas generativos, las plataformas de producción y consumo, la conversión de recipes en activos y la confianza como moneda requieren validarse con productos y mercados en los próximos años.

\begin{itemize}
\item ¿En qué tipo de contenido sustituirán primero los sistemas generativos a los sistemas de recomendación?
\item ¿Las plataformas se transformarán por sí mismas en plataformas de producción y consumo, o aparecerán primero plataformas nuevas?
\item ¿Cómo se pueden proteger los datos generados dentro del producto para evitar que los absorban las empresas de modelos base?
\item ¿Pueden la recipe, la IP y la confianza convertirse en activos con precio?
\item ¿Cómo se define la titularidad de los derechos de autor del contenido generado por IA en escenarios donde «consumir es crear»?
\item ¿Puede un estudio pequeño que trabaja a distancia sostener a largo plazo un sistema generativo complejo?
\end{itemize}

\subsection{Lecturas adicionales}

\begin{itemize}
\item EP130, entrevista con Zhang Yueguang (张月光): del diseño de procesos al diseño de contexto, productos AI Native y Agent.
\item EP128, entrevista con Manus: conversión de Agent en producto, riesgos de la complejidad y la IA como manufactura.
\item EP134, panorama de datos: pirámide de datos, Data Recipe y fijación de precios de los datos.
\item EP139, panorama de Agent: historia técnica de los Agent, OpenClaw Moment e impacto social.
\item Let's Verify Step by Step, Language Modeling Is Compression y otros materiales: para entender la supervisión de procesos, la compresión y la generalización.
\end{itemize}

\begin{importantbox}{Juicio final}
Lo que más vale la pena conservar de EP123 es una nueva visión de plataforma: cuando el sistema de producción es lo bastante potente, la plataforma ya no solo asigna inventario, sino que produce contenido según la intención del usuario. La moneda de la atención del internet anterior no desaparecerá, pero los activos de mayor valor se desplazarán hacia la recipe, la IP, el gusto y la confianza.
\end{importantbox}

\end{document}
